% =============================================================================
%  Appendix: Theory (Prop. 1-2, Thm. 3, Thm. 4 easy part)
%  To \input{} into the CVPR supplementary. Required in the preamble:
%    \usepackage{amsmath,amssymb,amsthm}
%    \newtheorem{theorem}{Theorem}  \newtheorem{proposition}[theorem]{Proposition}
%    \newtheorem{lemma}[theorem]{Lemma} \newtheorem{corollary}[theorem]{Corollary}
%    \theoremstyle{definition} \newtheorem{assumption}{Assumption}
%    \theoremstyle{remark} \newtheorem{remark}[theorem]{Remark} \newtheorem{example}[theorem]{Example}
%  Lines marked  % CHECK:  are things to verify against the cited source before submission.
% =============================================================================

\providecommand{\E}{\mathbb{E}}
\providecommand{\R}{\mathbb{R}}
\providecommand{\Prob}{\mathbb{P}}
\providecommand{\Law}{\mathrm{Law}}
\providecommand{\Unif}{\mathrm{Unif}}
\providecommand{\Leb}{\mathcal{L}}
\providecommand{\indep}{\perp\!\!\!\perp}
\newcommand{\Xs}{X^{*}}
\newcommand{\Ds}{D^{*}}
\newcommand{\Adm}{\mathcal{A}}
\newcommand{\Wtwo}{W_2}

\section{Theory}
\label{sec:app-theory}

\subsection{Setting and notation}
\label{sec:app-setting}

Let $X\in\R^d$ be the clean latent with law $p_X$ and $\E\|X\|^2<\infty$.
The encoder computes $y=g_a(X)\in\R^k$ (any side information, \eg a hyperprior, can be appended to $y$; nothing below uses the form of $g_a$).
For a step $\Delta>0$ let $C=[-\Delta/2,\Delta/2)$, let $Q_\Delta(v)=\Delta\lfloor v/\Delta+1/2\rfloor$ act componentwise, and let $r(v)=v-Q_\Delta(v)\in C^k$ be the quantization residual.
Encoder and decoder share a dither $U\sim\Unif(C^k)$ independent of $X$ (a shared seed).
The encoder transmits $Q=Q_\Delta(y+U)$ and the decoder forms
\begin{equation}
  \hat y = Q-U = y+E, \qquad E := Q_\Delta(y+U)-(y+U) .
\end{equation}
The decoder's information is $Y:=(Q,U)$. A \emph{decoder} (estimator) is any random vector $\hat X$ with $\E\|\hat X\|^2<\infty$ such that $X - Y - \hat X$ is a Markov chain, \ie $\hat X \indep X \mid Y$; this allows private decoder randomness. We write $\Adm$ for this set, and
\begin{equation}
  \Xs=\E[X\mid Y],\qquad \Ds=\E\|X-\Xs\|^2,\qquad W=\Wtwo(p_{\Xs},p_X),
\end{equation}
where $\Wtwo$ is the quadratic Wasserstein distance. Since $\E\|\Xs\|^2\le\E\|X\|^2<\infty$, $W<\infty$. The distortion--perception function is
\begin{equation}
  D(P)=\inf\bigl\{\E\|\hat X-X\|^2 : \hat X\in\Adm,\ \Wtwo(p_{\hat X},p_X)\le P\bigr\},\qquad P\ge 0 .
\end{equation}

\begin{lemma}[Orthogonality]\label{lem:orth}
For every $\hat X\in\Adm$, \quad $\E\|\hat X-X\|^2=\Ds+\E\|\hat X-\Xs\|^2$.
\end{lemma}
\begin{proof}
Expand $\hat X-X=(\hat X-\Xs)+(\Xs-X)$. The cross term is integrable by Cauchy--Schwarz, and by conditional independence of $\hat X$ and $X$ given $Y$,
$\E\langle \hat X-\Xs,\Xs-X\rangle=\E\bigl\langle \E[\hat X\mid Y]-\Xs,\ \Xs-\E[X\mid Y]\bigr\rangle=0$,
since $\Xs$ is $\sigma(Y)$-measurable and $\E[X\mid Y]=\Xs$.
\end{proof}

\subsection{Proposition 1: subtractive dither}

\begin{proposition}[Dither gives an exact, continuous source]\label{prop:dither}
In the setting above:
\begin{enumerate}
  \item[(a)] $E\sim\Unif(C^k)$ and $E\indep X$ (indeed $E\indep(X,y)$), hence $\hat y=y+E$ is $y$ corrupted by known, signal-independent uniform noise.
  \item[(b)] $\sigma(Y)=\sigma(\hat y)$ up to null sets; thus $\Xs=m(\hat y)$ with $m(v)=\E[X\mid\hat y=v]$.
  \item[(c)] $\hat y$ has a Lebesgue density on $\R^k$, $p_{\hat y}(v)=\Delta^{-k}\,\Prob(v-y\in C^k)\le\Delta^{-k}$, whatever the law of $y$ (even if $y$ is discrete).
  \item[(d)] If $k=d$ and $m$ has a version that is locally Lipschitz with $\det Dm\neq 0$ Lebesgue-a.e., then $p_{\Xs}\ll\Leb^d$.
  \item[(e)] Without dither ($U\equiv 0$), $Y=Q_\Delta(y)\in\Delta\mathbb{Z}^k$ is countably valued, so $p_{\Xs}$ is purely atomic. If $p_X$ has no atoms, there is no measurable $T$ with $T_\#p_{\Xs}=p_X$: every \emph{deterministic} decoder has $\Wtwo(p_{\hat X},p_X)>0$.
\end{enumerate}
\end{proposition}

\begin{proof}
(a) Condition on $X=x$; then $y=g_a(x)$ is fixed and $U$ is still $\Unif(C^k)$. Coordinates are independent, so fix one coordinate and a scalar $y$. Then $V=y+U\sim\Unif[y-\Delta/2,y+\Delta/2)$ and $E=-r(V)$. The residual $r$ is $\Delta$-periodic, and on any half-open interval of length $\Delta$ it is a piecewise translation onto $C$ (at most two pieces), hence it maps the uniform law on that interval to $\Unif(C)$. So $r(V)\sim\Unif(C)$, and $E=-r(V)$ is uniform on $-C$, which equals $\Unif(C)$ up to the null boundary. The conditional law of $E$ given $X=x$ does not depend on $x$, which proves independence. This is the classical result of \cite{schuchman1964dither,gray1993dithered}.

(b) Since $Q\in\Delta\mathbb{Z}^k$ and $-U\in -C$, we have $r(\hat y)=r(-U)=-U$ except on the null event $\{U_i=-\Delta/2 \text{ for some } i\}$. So $U=-r(\hat y)$ and $Q=\hat y+U$ a.s. Conversely $\hat y=Q-U$.

(c) For a Lebesgue-null Borel set $N\subset\R^k$, independence in (a) gives $\Prob(\hat y\in N)=\E\bigl[\Delta^{-k}\Leb^k\bigl((N-y)\cap C^k\bigr)\bigr]=0$. The same computation with $N$ a general Borel set yields the stated density.

(d) Let $N\subset\R^d$ be Lebesgue-null. By the area formula for locally Lipschitz maps \cite[Thm.~3.9]{evans2015measure},
$\int_{m^{-1}(N)}|\det Dm(v)|\,dv=\int_N \#\bigl(m^{-1}(z)\bigr)\,dz=0$.
So $|\det Dm|=0$ a.e.\ on $m^{-1}(N)$, and since $\det Dm\neq0$ a.e., $m^{-1}(N)$ is null. By (c), $\Prob(\Xs\in N)=\Prob(\hat y\in m^{-1}(N))=0$.
% CHECK: theorem number of the area formula in the edition of Evans & Gariepy you cite.

(e) $\Xs$ is a function of the countably valued $Y$, so $p_{\Xs}$ is purely atomic, and so is $T_\#p_{\Xs}$ for any measurable $T$. A purely atomic measure cannot equal an atomless $p_X$.
\end{proof}

\begin{remark}[Dimension reduction]\label{rem:dimred}
In latent codecs usually $k<d$ (the analysis transform reduces dimension). Then $\Xs=m(\hat y)$ lies in $m(\R^k)$, which is $\Leb^d$-null when $m$ is locally Lipschitz. So $p_{\Xs}$ is \emph{not} absolutely continuous and a Brenier map need not exist. Theorem~\ref{thm:ot-bridge}(a,b) below does not need absolute continuity: it uses an optimal \emph{plan} and decoder-side randomness. This is exactly the role of the start-point noise $\sigma(t_0)\varepsilon$ in our decoder. Part (c) of the theorem, the deterministic map, applies when $k=d$ (as in the toy experiment) or, heuristically, when $p_{\Xs}$ is ``nearly'' absolutely continuous.
\end{remark}

\subsection{Proposition 2: a single Euler step returns a conditional mean}

Let $(X_0,X_1)$ be any coupling with finite second moments, let $X_t=(1-t)X_0+tX_1$, and let $v(x,t)=\E[X_1-X_0\mid X_t=x]$ be the minimizer of the population flow-matching objective. It is defined $p_{X_t}$-a.e.\ for a.e.\ $t$; we use the version at $t=0$, which is the limit a velocity network that is continuous in $t$ attains.

\begin{proposition}[Euler step from the MMSE start]\label{prop:euler}
\leavevmode
\begin{enumerate}
  \item[(a)] $v(x,0)=\E[X_1\mid X_0=x]-x$. Hence an Euler step of size $h\in[0,1]$ from $t=0$ gives $X_0+h\,v(X_0,0)=(1-h)X_0+h\,\E[X_1\mid X_0]$.
  \item[(b)] For the natural coupling $(X_0,X_1)=(\Xs,X)$: $v(\cdot,0)=0$ $p_{\Xs}$-a.e. One Euler step returns $\Xs$ unchanged.
  \item[(c)] For the independent coupling $X_0\indep X_1$: one full Euler step returns the constant $\E[X]$.
  \item[(d)] Let $X_1\sim p_X$ and $\hat X=\E[X_1\mid X_0]$ (one full Euler step). Then $\E\|X_1\|^2=\E\|\hat X\|^2+\E\|X_1-\hat X\|^2$, and $\Wtwo(p_{\hat X},p_X)=0$ if and only if $X_1=f(X_0)$ a.s.\ for some measurable $f$, \ie the coupling is deterministic.
\end{enumerate}
\end{proposition}

\begin{proof}
(a) At $t=0$, $X_t=X_0$, so $v(x,0)=\E[X_1-X_0\mid X_0=x]=\E[X_1\mid X_0=x]-x$.
(b) $\sigma(\Xs)\subseteq\sigma(Y)$, so by the tower property $\E[X\mid\Xs]=\E[\E[X\mid Y]\mid\Xs]=\E[\Xs\mid\Xs]=\Xs$. Apply (a).
(c) $\E[X_1\mid X_0]=\E[X_1]=\E[X]$.
(d) The identity is the Pythagorean property of conditional expectation. If $\hat X\sim p_X$, then $\E\|\hat X\|^2=\E\|X_1\|^2$, so $\E\|X_1-\hat X\|^2=0$ and $X_1=\hat X$ is a function of $X_0$. Conversely, if $X_1=f(X_0)$ then $\hat X=f(X_0)=X_1\sim p_X$.
\end{proof}

\noindent\emph{Consequence.} A one-step Euler decoder started at the MMSE estimate is either inert (natural coupling), or outputs a conditional mean and therefore loses second moment, which shows up as blur (any non-deterministic coupling). It reaches perfect perception only if the training coupling is a deterministic map. Theorem~\ref{thm:ot-bridge} identifies which map is optimal.

\subsection{Theorem 3: the OT bridge traces the D--P frontier}

\begin{theorem}[OT bridge]\label{thm:ot-bridge}
Let $\gamma^*$ be an optimal coupling for $\Wtwo(p_{\Xs},p_X)$.
\begin{enumerate}
  \item[(a)] \textbf{D--P function} \cite{freirich2021theory}. For all $P\ge0$, \ $D(P)=\Ds+\bigl((W-P)_+\bigr)^2$.
  \item[(b)] \textbf{Achievability.} There is $Z$ with $(\Xs,Z)\sim\gamma^*$ and $\hat x_t:=(1-t)\Xs+tZ\in\Adm$ for all $t\in[0,1]$, and
  \begin{equation}\label{eq:geodesic}
    \E\|\hat x_t-X\|^2=\Ds+t^2W^2,\qquad \Wtwo(p_{\hat x_t},p_X)=(1-t)\,W .
  \end{equation}
  So $\hat x_t$ attains $D(P)$ at $P=(1-t)W$, and $t\in[0,1]$ parametrizes the entire frontier.
  \item[(c)] \textbf{Deterministic decoder.} If $p_{\Xs}\ll\Leb^d$ (Prop.~\ref{prop:dither}(d)), then $\gamma^*=(\mathrm{id},T)_\#p_{\Xs}$ with $T=\nabla\varphi$ the Brenier map \cite{brenier1991polar}, $Z=T(\Xs)$, and $\hat x_t=(1-t)\Xs+t\,T(\Xs)$ is a deterministic function of $Y$.
  \item[(d)] \textbf{Flow matching with the OT coupling.} Under (c), train flow matching on $(X_0,X_1)=(\Xs,T(\Xs))$. Then $v(x,0)=T(x)-x$, so \emph{one Euler step of size $t$ from $t=0$ outputs exactly $\hat x_t$}. Moreover, for every $t\in[0,1)$ the map $x\mapsto(1-t)x+tT(x)$ is injective on the set where $T$ is defined. Hence $v(x,t)=(T-\mathrm{id})\circ S_t^{-1}(x)$ with $S_t=(1-t)\mathrm{id}+tT$, the straight lines $t\mapsto(1-t)x_0+tT(x_0)$ solve $\dot x=v(x,t)$, and the straightness gap $\E\int_0^1\|v(X_t,t)-(X_1-X_0)\|^2dt$ is zero.
\end{enumerate}
\end{theorem}

\begin{proof}
\emph{Step 0 (existence of $Z$).} $p_{\Xs}$ and $p_X$ have finite second moments, so an optimal coupling $\gamma^*$ exists \cite[Thm.~4.1]{villani2009optimal}. Disintegrate $\gamma^*(dx,dz)=p_{\Xs}(dx)\,\kappa(x,dz)$. Since $\R^d$ is Polish, there is a measurable $G$ such that $G(x,\xi)\sim\kappa(x,\cdot)$ for $\xi\sim\Unif[0,1]$ \cite[Lemma~4.22]{kallenberg2021foundations}. Take $\xi$ independent of $(X,Y)$ and set $Z=G(\Xs,\xi)$. Then $(\Xs,Z)\sim\gamma^*$, and $\hat x_t$ is a function of $(Y,\xi)$ with $\xi\indep(X,Y)$. So $X-Y-\hat x_t$ is Markov and $\hat x_t\in\Adm$.
% CHECK: lemma number of the transfer/randomization lemma in the Kallenberg edition you cite.

\emph{Proof of (b).} By Lemma~\ref{lem:orth}, $\E\|\hat x_t-X\|^2=\Ds+\E\|\hat x_t-\Xs\|^2=\Ds+t^2\E\|Z-\Xs\|^2=\Ds+t^2W^2$, since $\gamma^*$ is optimal. For the perception term, the coupling $(\hat x_t,Z)$ gives $\Wtwo(p_{\hat x_t},p_X)\le(\E\|\hat x_t-Z\|^2)^{1/2}=(1-t)W$. Likewise $(\Xs,\hat x_t)$ gives $\Wtwo(p_{\Xs},p_{\hat x_t})\le tW$. The triangle inequality for $\Wtwo$ then gives $W\le\Wtwo(p_{\Xs},p_{\hat x_t})+\Wtwo(p_{\hat x_t},p_X)\le tW+\Wtwo(p_{\hat x_t},p_X)$, \ie $\Wtwo(p_{\hat x_t},p_X)\ge(1-t)W$.

\emph{Proof of (a).} \emph{Lower bound:} let $\hat X\in\Adm$ with $\Wtwo(p_{\hat X},p_X)\le P$. Since $(\hat X,\Xs)$ is a coupling, Lemma~\ref{lem:orth} gives $\E\|\hat X-X\|^2=\Ds+\E\|\hat X-\Xs\|^2\ge\Ds+\Wtwo(p_{\hat X},p_{\Xs})^2$. The triangle inequality gives $\Wtwo(p_{\hat X},p_{\Xs})\ge W-\Wtwo(p_{\hat X},p_X)\ge W-P$. Since distances are nonnegative, $\Wtwo(p_{\hat X},p_{\Xs})\ge(W-P)_+$, and squaring gives $\E\|\hat X-X\|^2\ge\Ds+((W-P)_+)^2$. \emph{Upper bound:} if $P\ge W$, $\hat X=\Xs\in\Adm$ achieves $\Ds$. If $0\le P<W$, $\hat x_t$ with $t=1-P/W$ achieves $\Ds+(W-P)^2$ by (b).

\emph{Proof of (c).} Brenier's theorem \cite{brenier1991polar}; see also \cite[Thm.~2.12]{villani2003topics}.

\emph{Proof of (d).} Here $X_1=T(X_0)$, so Prop.~\ref{prop:euler}(a) gives $v(x,0)=T(x)-x$, and $X_0+t\,v(X_0,0)=\hat x_t$. For injectivity, $T=\nabla\varphi$ with $\varphi$ convex, so on the set where $\nabla\varphi$ exists, monotonicity of the gradient gives
$\langle S_t(x)-S_t(x'),x-x'\rangle=(1-t)\|x-x'\|^2+t\langle\nabla\varphi(x)-\nabla\varphi(x'),x-x'\rangle\ge(1-t)\|x-x'\|^2$.
So $S_t$ is injective for $t<1$ and $X_0=S_t^{-1}(X_t)$ is a function of $X_t$. Then $v(X_t,t)=\E[T(X_0)-X_0\mid X_t]=(T-\mathrm{id})(S_t^{-1}(X_t))=X_1-X_0$ a.s., which is the zero-gap claim. Along $x_t=S_t(x_0)$ we have $\dot x_t=T(x_0)-x_0=v(x_t,t)$.
\end{proof}

\begin{corollary}[Why the OT coupling]\label{cor:why-ot}
Let $f$ be measurable with $f_\#p_{\Xs}=p_X$ (any deterministic coupling, hence an exact one-step decoder by Prop.~\ref{prop:euler}). Then $\hat X=f(\Xs)$ has perfect perception and
$\E\|f(\Xs)-X\|^2=\Ds+\E\|f(\Xs)-\Xs\|^2\ge\Ds+W^2$,
with equality if and only if $(\mathrm{id},f)_\#p_{\Xs}$ is an optimal coupling, \ie $f=T$ $p_{\Xs}$-a.e.\ under (c).
\end{corollary}
\begin{proof}
Lemma~\ref{lem:orth} and the definition of $W$. Uniqueness of the optimal coupling under (c) is part of Brenier's theorem.
\end{proof}

\begin{remark}[What is and is not claimed]\label{rem:straight}
Theorem~\ref{thm:ot-bridge}(d) and Corollary~\ref{cor:why-ot} give: OT coupling $\Rightarrow$ zero straightness gap and D--P optimality. The converse (``straight $\Rightarrow$ OT'') fails in general for $d\ge2$. Straight, non-crossing couplings, \eg fixed points of rectification, need not be optimal for the quadratic cost \cite{liu2023flow}.
% CHECK: cite the exact place in Liu et al. (2023) that discusses straight vs. c-optimal couplings.
Also, one-step exactness holds for \emph{any} deterministic coupling (Prop.~\ref{prop:euler}(d)). What singles out OT is the distortion (Corollary~\ref{cor:why-ot}).
\end{remark}

\subsection{Theorem 4 (first part): the rectified decoder is within $2\times$ MMSE}

\begin{assumption}\label{ass:rf}
For the natural coupling $(X_0,X_1)=(\Xs,X)$, the marginal velocity $v$ is such that $\dot z=v(z,t)$ has a unique solution on $[0,1]$ from $p_{\Xs}$-a.e.\ start. Its flow map $\Phi$ satisfies $\Law(\Phi_t(\Xs))=\Law(X_t)$ for all $t$ (marginal preservation).
% CHECK: match these to the exact hypotheses of the marginal-preservation theorem in Liu et al. (2023).
\end{assumption}

\begin{theorem}[Rectified decoder]\label{thm:rectified}
Under Assumption~\ref{ass:rf}, let $\hat X=\Phi_1(\Xs)$. Then $\hat X\in\Adm$, $\Wtwo(p_{\hat X},p_X)=0$, and
\begin{equation}
  \Ds+W^2\ \le\ \E\|\hat X-X\|^2\ \le\ 2\Ds .
\end{equation}
The upper bound equals the distortion of posterior sampling $\hat X\sim p_{X\mid Y}$, which also attains perfect perception. The rectified decoder attains it deterministically.
\end{theorem}
\begin{proof}
$\hat X$ is a function of $\Xs$, hence of $Y$, so $\hat X\in\Adm$. Marginal preservation at $t=1$ gives $\hat X\sim p_X$. By the convex-cost reduction property of rectification \cite{liu2023flow} with $c(z)=\|z\|^2$,
$\E\|\hat X-\Xs\|^2\le\E\|X-\Xs\|^2=\Ds$.
% CHECK: theorem number (convex transport cost does not increase under rectification) in Liu et al. (2023).
Lemma~\ref{lem:orth} gives $\E\|\hat X-X\|^2=\Ds+\E\|\hat X-\Xs\|^2\le2\Ds$. The lower bound is Theorem~\ref{thm:ot-bridge}(a) at $P=0$. For posterior sampling, $\hat X\indep X\mid Y$ with the same conditional law, so $\E\|\hat X-X\|^2=2\,\E\,\mathrm{tr}\,\mathrm{Cov}(X\mid Y)=2\Ds$.
\end{proof}

By Prop.~\ref{prop:euler}(b) the rectified trajectory starts with zero velocity, $v(\cdot,0)=0$. All motion happens at $t>0$, which is why its one-step Euler discretization fails.

\subsection{Example: Gaussian source, dithered scalar quantization}
\label{sec:app-toy}

Let $X\sim\mathcal N(0,\mathrm{diag}(\sigma_1^2,\dots,\sigma_d^2))$ be quantized in its KLT basis ($k=d$, $g_a=\mathrm{id}$). Then $\hat y=X+E$ with $E\sim\Unif(C^d)$, and everything factorizes over coordinates. With $a_i=(\hat y_i-\Delta/2)/\sigma_i$ and $b_i=(\hat y_i+\Delta/2)/\sigma_i$:
\begin{align}
  \Xs_i&=m_i(\hat y_i)=\sigma_i\,\frac{\phi(a_i)-\phi(b_i)}{\Phi(b_i)-\Phi(a_i)}, \\
  T_i(\Xs_i)&=\sigma_i\,\Phi^{-1}\!\bigl(F_{\hat y_i}(\hat y_i)\bigr),\qquad
  F_{\hat y_i}(v)=\frac{\sigma_i}{\Delta}\Bigl[G\bigl(\tfrac{v+\Delta/2}{\sigma_i}\bigr)-G\bigl(\tfrac{v-\Delta/2}{\sigma_i}\bigr)\Bigr],
\end{align}
where $G(z)=z\Phi(z)+\phi(z)$. Each $m_i$ is smooth and strictly increasing, so Prop.~\ref{prop:dither}(d) applies. For product measures the coordinatewise monotone map is the gradient of the separable convex function $\sum_i\varphi_i(x_i)$ and pushes $p_{\Xs}$ to $p_X$, so it is the Brenier map. Hence $W^2=\sum_i\Wtwo^2(p_{\Xs_i},p_{X_i})$, and $\Ds$ and $W$ reduce to one-dimensional integrals. Note that $p_{\Xs}$ is not Gaussian here, so the Gaussian (Bures) closed form does \emph{not} apply.
